\documentclass[main.tex]{subfiles}

\begin{document}

\section{Objectifs et contraintes}

\subsection{Les objectifs techniques}

\quad Le projet Casiers RFID Motorisés repose sur plusieurs défis techniques visant à garantir son efficacité, sa fiabilité et sa simplicité d’utilisation. Les principaux objectifs techniques sont les suivants :

\begin{enumerate}

	\item Automatisation de l’ouverture des casiers :
	\begin{itemize}
		\item Intégration de servomoteurs intelligents pour l’ouverture et la fermeture des casiers
		\item Indication visuelle avec des diodes RGB pour signaler le casier concerné
	\end{itemize}
	\item Gestion des identifiants par RFID :
	\begin{itemize}
		\item Utilisation de badges RFID UM4100 à 125 kHz pour identifier les chauffeurs
		\item Développement d’une application Windows permettant l’enregistrement et la lecture des badges
	\end{itemize}
	\item Base de données et traçabilité :
	\begin{itemize}
		\item Conception d’une base de données SQL hébergée sur une machine virtuelle sous Proxmox
		\item Enregistrement des mouvements des chauffeurs pour une exploitation statistique ultérieure
	\end{itemize}
	\item Sécurisation et communication réseau :
	\begin{itemize}
		\item Intégration d’un pare-feu logiciel (IPFire) pour protéger le système
		\item Utilisation du réseau Ethernet interne de l’usine pour la synchronisation des données
	\end{itemize}
	
\end{enumerate}

L’ensemble de ces aspects doit être intégré de manière cohérente et optimisée, afin de garantir un fonctionnement fiable, rapide et sécurisé du système.


\subsection{Les objectifs économiques}

\quad Le projet étant développé pour l’entreprise Sylvamo dans le cadre d’une optimisation des flux logistiques, il doit répondre à plusieurs contraintes économiques :

\begin{enumerate}

	\item Réduction des coûts de gestion des entrées et sorties :
	\begin{itemize}
		\item Diminution du temps passé par les agents de sécurité sur la gestion des documents
		\item Fluidification du trafic pour éviter des pertes de temps coûteuses pour les transporteurs
	\end{itemize}
	\item Coût maîtrisé du matériel :
	\begin{itemize}
		\item Budget alloué de 1000 € pour l’ensemble du projet
		\item Utilisation de composants économiques et faciles à approvisionner (Raspberry Pi, servomoteurs abordables, badges RFID standards)
	\end{itemize}
	\item Durabilité et maintenance minimale :
	\begin{itemize}
		\item Choix de matériaux et composants garantissant une longue durée de vie
		\item Minimisation des interventions humaines pour réduire les coûts de maintenance
	\end{itemize}

\end{enumerate}

Le système doit donc être rentable sur le long terme en améliorant l’efficacité du processus sans engendrer des coûts excessifs.

\newpage

\subsection{Les délais}

\quad Le projet s’inscrit dans un planning strict qui doit être respecté pour assurer sa finalisation dans le cadre du BTS CIEL. Voici les principales échéances :

\begin{itemize}

	\item Janvier 2025 : Lancement du projet et début de l’analyse des besoins
	\item Février - Mars 2025 :
	\begin{itemize}
		\item Conception et modélisation (diagrammes UML, architecture du système)
		\item Développement de la base de données et des premières fonctionnalités logicielles
	\end{itemize}
	\item Avril 2025 :
	\begin{itemize}
		\item Assemblage du prototype physique et intégration des composants électroniques
		\item Tests unitaires des différentes fonctionnalités
	\end{itemize}
	\item Mai 2025 :
	\begin{itemize}
		\item Finalisation du projet et tests en condition réelle
		\item Rédaction du rapport et préparation de la soutenance
	\end{itemize}
	\item Juin 2025 
	\begin{itemize}
		\item Présentation du projet
	\end{itemize}
	
\end{itemize}

\newpage

\end{document}
